\documentclass[11pt]{article}
\usepackage[margin=1in]{geometry}
\usepackage{amsmath}
\usepackage{enumitem}
\usepackage{xcolor}
\usepackage{hyperref}
\setlist[description]{leftmargin=1.4em,style=nextline,font=\normalfont\itshape}
\newcommand{\concern}[1]{\par\smallskip\textbf{Concern.} #1}
\newcommand{\response}[1]{\par\textbf{Response.} #1}
\newcommand{\action}[1]{\par\textbf{Action.} #1\par\smallskip}
\newcommand{\nmpc}{NMPC}
\title{\vspace{-2.5em}Response to Reviewers}
\author{}
\date{}
\begin{document}
\maketitle
\vspace{-3em}

\noindent\textbf{Manuscript \#:} Access-2026-22287\\
\textbf{Title:} Point Mass Planning and \nmpc{} Tracking for Time-Optimal
Agile Quadrotor Gate Navigation\\
\textbf{Authors:} B.\ S.\ Guevara, A.\ S.\ Brand\~ao, J.\ Varela-Ald\'as

\medskip
\noindent We thank the Associate Editor and both reviewers. Both found the
work technically sound and a contribution to the field, and we have
addressed every comment. Changes are marked in yellow in the highlighted
PDF; section/equation labels refer to the revised manuscript. No new
references were added.

\section*{Reviewer 1}
The reviewer notes we identified and addressed two significant research
gaps, and answered \emph{Yes} to contribution, technical soundness, and
reference adequacy. We address the two suggestions below.

\subsection*{Comment 1.1 --- more detailed explanation of the concepts (Q2)}
\concern{The planner/controller description, while sufficient, could be
more detailed.}
\response{We added insight at the one point previously stated without
motivation---the point-mass cost structure---as it links the planner to the
quality of the reference passed to the tracker.}
\action{A sentence after Eq.~(PMM cost, Sec.~III-C) now explains the three
cost terms, in particular that the jerk regulariser ($w_j$) keeps the
downstream flatness derivatives (attitude, body rate) well conditioned,
since they are obtained by differentiating the planned acceleration.}

\subsection*{Comment 1.2 --- criterion for the \nmpc{} parameters $N$, $f_c$ (Q3)}
\concern{No criterion is given for selecting the \nmpc{} parameters.}
\response{The values follow standard practice and the real-time budget, not
a separate tuning study, which we now state. $f_c=100$~Hz is the
conventional rate for body-rate agile \nmpc{}; $N=50$ ($T_h=0.5$~s) is the
largest horizon keeping the P99 solve time within the 10~ms deadline on the
test platform, and the half-second look-ahead spans roughly one inter-gate
transit ($T_f^*/N_g\approx0.7$~s).}
\action{Added an explanatory passage in Sec.~V-A after the parameter list,
giving the criterion for $f_c$ and $N$/$T_h$ and linking it to the measured
real-time budget (Sec.~V-G).}

\subsection*{Remaining items}
Questions 1, 4, and 5 were affirmative and required no action. No references
were suggested for addition or removal.

\section*{Reviewer 2}
The reviewer found the paper ``very well written, well-founded and
referenced'' and recommended publication after minor adjustments.

\subsection*{Comment 2.1 --- discuss the contribution against similar cited works}
\concern{Include a brief discussion positioning the proposal against
similar cited works, highlighting its advantages.}
\response{We agree; the paper did not consolidate this comparison in one
place. We added it using only references already in the manuscript.}
\action{A paragraph at the end of Sec.~V-H (Discussion) now positions the
work: it retains the fast point-mass planning of Foehn~et~al.\ but feeds the
\emph{complete} flat output to the \nmpc{}; it keeps planner and tracker
modular---unlike the fused MPCC of Romero~et~al.\ and Krinner~et~al.,
whose stages cannot be replaced independently---while recovering the
tracking completeness those methods achieve implicitly; and, unlike
Sun~et~al., it \emph{ablates} the reference structure to isolate which flat
components drive gate-crossing performance, at no extra solver cost.}

\subsection*{Remaining items}
All four assessment questions were affirmative; no references were suggested
for addition or removal.

\section*{Additional corrections by the authors}
We audited every quantitative claim against the raw result files. Tables~III
and~IV, all reported percentages, peak speeds, and peak accelerations were
confirmed correct. Two items were corrected:

\begin{description}
\item[AC1 --- Blow-up exclusion.] The text gave an approximate per-mode
divergence rate. We replaced it with the exact paired-exclusion outcome:
trials are dropped if either controller diverges ($\text{RMSE}>5$~m),
leaving 23/25 valid trials on the figure-8 and 18/25 on the vertical loop
(Sec.~V-D).
\item[AC2 --- Figure~2.] The original figure had been rendered from an
earlier planner export and the model-in-the-loop set, so its numbers did
not match the software-in-the-loop tables. It is regenerated from exactly
the reported data (Table~III reference; software-in-the-loop trajectories).
\end{description}

\medskip
\noindent We believe the revision fully addresses the reviewers' comments
and thank the reviewers and Associate Editor.

\end{document}
